\chapter{Chạy Xe Cũng Không Buông}
\markboth{Chạy Xe Cũng Không Buông}{Chạy Xe Cũng Không Buông}

\begin{center}
	\textit{\color{bookprimary}``Có những thói quen không chỉ vô duyên. Chúng có thể giết người.''}

	\textit{\color{bookprimary}``Some habits are not merely rude. They can kill.''}
\end{center}

\ngancach

\begin{tcolorbox}[colback=booksecondary!8,colframe=booksecondary!40,arc=5pt,title={\bfseries Góc Suy Nghĩ},fonttitle=\bfseries\color{booksecondary}]
Xem điện thoại khi chạy xe không còn là chuyện tập trung kém.
\textit{(Looking at a phone while driving is no longer just about poor concentration.)}

Nó là chuyện coi thường mạng sống.
\textit{(It is about disrespecting life.)}

Không chỉ mạng sống của mình.
\textit{(Not only your own life.)}

Mà còn của người đang đi cạnh, đi đối diện, hoặc vô tình đi ngang qua lúc em cúi xuống đúng hai giây ngu ngốc đó.
\textit{(Also the lives of those beside you, opposite you, or simply passing by during those two stupid seconds when you look down.)}
\end{tcolorbox}

\section{Một Tin Nhắn Không Đáng Mạng Người}
\begin{truyen}{Một Tin Nhắn Không Đáng Mạng Người}{No Message Is Worth a Life}
\chuhoa{C}{ó người vừa chạy xe vừa trả lời tin nhắn bằng một tay.}
\textit{(Someone rides a motorbike while replying to a message with one hand.)}

Khi bị nhắc, người đó nói rất quen: ``Chỉ vài giây thôi mà.''
\textit{(When warned, the familiar answer comes: ``It is only a few seconds.'')}

\adultpair{Đúng. Tai nạn lớn nhiều khi cũng chỉ cần vài giây. Nhiều người tưởng nguy hiểm là thứ phải ồn ào, rõ ràng, báo trước rất lâu. Không. Có những tai nạn chỉ cần một cú liếc xuống, một bàn tay lỏng lái, một chiếc xe từ hẻm lao ra, là mọi lời biện hộ xong đời luôn.}{``Exactly. Major accidents often need only a few seconds too. Many people think danger must be loud, obvious, and announced well in advance. No. Some accidents need only one downward glance, one loose grip on the handlebar, one bike emerging from an alley, and every excuse is finished forever.''}

Người lớn nào làm vậy trước mặt con trẻ đang dạy một bài rất độc.
\textit{(Any adult doing that in front of children is teaching a poisonous lesson.)}

Bài đó là: sự tiện của mình quan trọng hơn sự an toàn của tất cả.
\textit{(That lesson is: my convenience matters more than everyone else's safety.)}
\end{truyen}

\section{Đèn Đỏ Không Phải Lúc Trả Lời}
\begin{truyen}{Đèn Đỏ Không Phải Lúc Trả Lời}{A Red Light Is Not Reply Time}
\chuhoa{C}{ó người tự cho rằng dừng đèn đỏ vài chục giây thì tranh thủ xem điện thoại là vô hại.}
\textit{(Some people assume that checking the phone for a few seconds at a red light is harmless.)}

Nhưng thói quen xấu hiếm khi đứng yên đúng trong giới hạn mình tự vẽ.
\textit{(But bad habits rarely stay inside the limit we draw for them.)}

Từ đèn đỏ, nó tràn sang lúc xe vừa nhích, vừa cua, vừa chen.
\textit{(From the red light, it spills into the moment the vehicle starts moving, turns, or squeezes through traffic.)}

\adultpair{Người ta không chỉ tập một động tác. Người ta tập cả nếp sống. Khi em tập cho mình phản xạ đụng vào điện thoại giữa dòng giao thông, sớm muộn phản xạ đó cũng chen vào lúc nguy hiểm hơn.}{``People do not practice only one action. They practice a whole way of living. If you train yourself to touch the phone in traffic, that reflex will sooner or later intrude at a more dangerous moment.''}
\end{truyen}

\section{Cầm Điện Thoại Là Cầm Thêm Rủi Ro}
\begin{truyen}{Cầm Điện Thoại Là Cầm Thêm Rủi Ro}{Holding the Phone Means Holding Extra Risk}
\chuhoa{C}{ó người nghĩ mình lái xe quen nên một tay vẫn giữ tay lái được.}
\textit{(Some people think their driving experience is enough to handle the handlebars with one hand.)}

Kinh nghiệm không làm đường bớt bất ngờ.
\textit{(Experience does not make the road less unpredictable.)}

Nó cũng không làm phản xạ vật lý nhanh hơn định luật va chạm.
\textit{(Nor does it make physical reaction faster than collision.)}

\adultpair{Khi em cầm điện thoại lúc chạy xe, em không chỉ cầm một vật. Em cầm thêm xác suất gây chuyện cho mình và cho người khác. Mà xác suất đó không xin lỗi ai trước khi xảy ra.}{``When you hold a phone while driving, you are not holding only an object. You are also holding a larger probability of harming yourself and others. And probability never apologizes before it happens.''}
\end{truyen}

\section{Con Ngồi Sau Lưng Đang Học Theo}
\begin{truyen}{Con Ngồi Sau Lưng Đang Học Theo}{The Child Behind You Is Learning from It}
\chuhoa{C}{ó phụ huynh vừa chở con vừa liếc điện thoại, trong khi đứa trẻ ngồi sau nhìn thấy hết.}
\textit{(A parent drives while glancing at the phone, and the child sitting behind sees everything.)}

Bài học ở đó không nằm trong lời nói.
\textit{(The lesson there is not in words.)}

Nó nằm trong cảnh tượng được lặp lại đủ lần để thành chuẩn hành vi.
\textit{(It lies in the scene repeated enough times to become a behavioral standard.)}

\adultpair{Trẻ con học cái được sống trước mặt nó nhanh hơn học cái bị giảng vào tai nó. Người lớn lái xe với điện thoại là đang âm thầm cấp giấy phép đạo đức cho thói vô trách nhiệm ấy.}{``Children learn what is lived in front of them faster than what is lectured into their ears. An adult driving with a phone is quietly issuing moral permission for that same irresponsibility.''}
\end{truyen}

\section{Tai Nạn Không Báo Trước Bằng Chuông}
\begin{truyen}{Tai Nạn Không Báo Trước Bằng Chuông}{Accidents Do Not Ring First}
\chuhoa{C}{ó người bảo rằng mình chỉ nhìn khi đường đang vắng.}
\textit{(Someone says they look only when the road seems empty.)}

Nhưng tai nạn vốn không có bổn phận xuất hiện đúng lúc mình đã chuẩn bị xong.
\textit{(But accidents have no obligation to appear only when we are ready.)}

\adultpair{Nguy hiểm thật không phải lúc nào cũng ồn. Nhiều khi nó đến đúng lúc em chủ quan nhất, ở một chỗ em tưởng mình kiểm soát được. Bởi vậy sự an toàn không thể đặt trên niềm tin mù vào thói quen cá nhân.}{``Real danger is not always loud. Often it arrives exactly when you are most confident, in a place you think you control. That is why safety cannot rest on blind trust in personal habit.''}
\end{truyen}

\section{Mọi Lý Do Đều Nhỏ Trước Mạng Sống}
\begin{truyen}{Mọi Lý Do Đều Nhỏ Trước Mạng Sống}{Every Reason Is Small Before a Human Life}
\chuhoa{C}{ó người nêu đủ lý do: công việc gấp, người nhà gọi, khách nhắn, đường còn xa.}
\textit{(Someone lists many reasons: urgent work, a family call, a client message, a long road ahead.)}

Nghe cái nào cũng có vẻ hợp lý.
\textit{(Each one sounds reasonable.)}

Nhưng tất cả đều nhỏ lại nếu đặt cạnh một mạng người không có cơ hội làm lại.
\textit{(But all of them shrink when placed beside a human life that gets no second attempt.)}

\adultpair{Người trưởng thành không phải người có nhiều lý do biện hộ. Là người biết thứ gì phải dừng lại trước ranh giới an toàn.}{``A mature person is not one with many justifications. It is one who knows what must stop at the boundary of safety.''}
\end{truyen}

\section{Người Cẩn Thận Không Tự Thử Vận May}
\begin{truyen}{Người Cẩn Thận Không Tự Thử Vận May}{Careful People Do Not Gamble with Luck}
\chuhoa{C}{ó người nói rằng mình làm vậy bao nhiêu lần mà có sao đâu.}
\textit{(Someone says they have done it many times and nothing happened.)}

Câu đó không chứng minh họ đúng.
\textit{(That sentence does not prove they were right.)}

Nó chỉ chứng minh họ đã may nhiều lần.
\textit{(It proves only that they have been lucky many times.)}

\adultpair{Đừng biến vận may thành triết lý sống. Người khôn không lấy những lần thoát nạn làm bằng chứng rằng hành vi của mình an toàn.}{``Do not turn luck into a philosophy of life. Wise people do not use repeated escapes as proof that their behavior is safe.''}
\end{truyen}

\section{Dừng Lại Hai Phút Đỡ Hối Cả Đời}
\begin{truyen}{Dừng Lại Hai Phút Đỡ Hối Cả Đời}{Pause for Two Minutes to Avoid a Lifetime of Regret}
\chuhoa{C}{ó người tập cho mình một quy tắc rất đơn giản: cần xem điện thoại thì tấp vào hẳn, dừng hẳn, rồi mới nhìn.}
\textit{(Someone trains a very simple rule: if the phone must be checked, pull over fully, stop fully, then look.)}

Chỉ một quy tắc nhỏ nhưng cứu được nhiều thứ lớn.
\textit{(One small rule can protect many large things.)}

\adultpair{Kỷ luật giao thông không cần phải bắt đầu từ triết lý cao siêu. Nó bắt đầu từ việc mình chấp nhận chậm hơn một chút để không biến người khác thành giá phải trả cho sự hấp tấp của mình.}{``Traffic discipline does not need to begin with lofty philosophy. It begins when we accept being a little slower so that other people do not pay the price for our haste.''}
\end{truyen}

\begin{baihoc}
Không có thông báo nào khẩn đến mức đáng để đổi bằng mạng người.
\textit{(No notification is so urgent that it deserves to be traded for a human life.)}

Nếu cần xem, hãy dừng lại.
\textit{(If you need to check, stop first.)}

Nếu không dừng nổi, vấn đề không còn là điện thoại nữa.
\textit{(If you cannot stop, the problem is no longer the phone alone.)}

Nó là nhân cách tự chủ của mình đang có lỗ thủng.
\textit{(It means your self-control as a person has a hole in it.)}
\end{baihoc}

\begin{tuhocnhanh}
1. Em đã từng xem điện thoại khi chạy xe hoặc đi đường chưa? \textit{(Have you ever checked your phone while riding or moving in traffic?)}

2. Nếu có, em đang đánh đổi điều gì để lấy vài giây đó? \textit{(If so, what are you trading away for those few seconds?)}

3. Em có dám nhắc một người lớn làm vậy không? \textit{(Would you dare to warn an adult doing that?)}
\end{tuhocnhanh}

\begin{tongketchuong}
Có những thói quen không đáng được gọi là thói quen xấu nữa.
\textit{(Some habits no longer deserve to be called merely bad habits.)}

Chúng là thói quen vô trách nhiệm với mạng người.
\textit{(They are habits of irresponsibility toward human life.)}
\end{tongketchuong}

\ngancach